% -*- mode: LaTeX; TeX-PDF-mode: t; -*-
\input{@resources/tex-add-search-paths}  % allow latex to find custom stuff
\input{./.econtexRoot}\documentclass[\econtexRoot/BufferStockTheory]{subfiles}

\newcommand{\subname}{IndividualBufferStockStability}
%\input{\econtexRoot/@local/dir-paths}
\usepackage{econark-ifsubfile}
\usepackage{econark-xrsetup}
\externaldocument{BufferStockTheory}

\xrsetup{\econtexRoot/\texname}

\begin{document}

\hypertarget{Analysis of the Converged Consumption Function}{}
\section{Individual Buffer Stock Stability}\label{sec:individStability}



In this section we analyse two notions of stability which will be useful for studying either an individual or a population of individuals who behave according to the converged consumption rule.
Consider an individual who at time $t$ holds normalized market resources $\mNrm_{t}$, and market resources in levels $\mLvl_{t}$, and follows the converged decision function $\cFunc$.
The time-$t$ consumption for the consumer will be $\cNrm_{t} = \cFunc(\mNrm_{t})$ and normalized market resources in time $t+1$ will be a random variable $\mNrm_{t+1} = \RNrmByGRnd_{t+1}(\mNrm_{t} - \cFunc(\mNrm_{t})) + \tranShkAll_{t+1}$.\footnote{None of the arguments in either of the two prior sections depended on the assumption that the consumption functions had converged.
With more cumbersome notation, each derivation could have been replaced by the corresponding finite-horizon versions.
This strongly suggests that it should be possible to extend the circumstances under which the problem can be shown to define a contraction mapping to the union of the parameter values under which \{\RIC,\FHWC\} hold and \{\FVAC,\WRIC\} hold.
That extension is not necessary for our purposes here, so we leave it for future work.}


Our first notion of stability concerns the existence of a unique `buffer stock target' for the individual; we are interested in whether the current level of normalized market resources is above or below a `target' level such that the magnitude of the precautionary motive (which causes a consumer to save) exactly balances the impatience motive  (which makes them want to dissave).
At the individual target, expected normalized market resources in the next period, \textit{conditioned on current normalized market resources}, will be the same as in the current period.
The intensifying strength of the precautionary motive with decreasing market resources can ensure stability of the target.
Below the target, the urgency to save due to the precautionary motive leads to an expected rise in market resources.
Conversely, above the target, impatience prevails, leading to an expected reduction of market resources.
%In this way, the `target' essentially defines the desired `buffer stock' of resources for the consumer.


Our second, weaker, notion of stability gives conditions for the invidiual such that an aggregate balanced growth path exists. To motivate this notion, consider Figure~\ref{fig:cNrmTargetFig} which shows the expected growth factors for consumption, the level of market resources, and normalized market resources, $\Ex_{t}[\cLvl_{t+1}/\cLvl_{t}]$, $\Ex_{t}[\mLvl_{t+1}/\mLvl_{t}]$, and $\Ex_{t}[\mNrm_{t+1}/\mNrm_{t}]$.\footnote{The figure is generated using parameters discussed in Section \ref{sec:GICdiscussion}, Table \ref{table:Calibration}.}
Begin by noting how the figure shows that as $\mNrm_{t} \rightarrow \infty$ the expected consumption growth factor goes to ${\APFac}$, indicated by the lower bound in Figure~\ref{fig:cNrmTargetFig}.
Moreover, as $\mNrm_{t}$ approaches zero, the consumption growth factor approaches $\infty$.
(Proposition  \ref{prop:convgGrowth}  in Appendix \ref{subsec:AppxCgrowthFac} establishes the asymptotic growth factors formally.)

Next, consider the implications of Figure~\ref{fig:cNrmTargetFig} for individual stability.
The figure shows a buffer stock target for normalized market resources, $\mNrm_{t} = \mBalLvl$, at which the expected growth factor of the level of market resources $\mLvl$ matches the expected growth factor of permanent income $\PermGroFac$.
A distinct and larger target ratio, $\mTrgNrm$, also exists.
At this ratio, $\Ex_{t}[\mNrm_{t+1}/\mNrm_{t}]=1$, and the expected growth factor of consumption is less than $\PermGroFac$.
Importantly, conditioned on an individual's time $t$ state, this model does not have a single $\mNrm$ at which $\permLvl$, $\mLvl$ and $\cLvl$ are all expected to grow at the same rate.
Yet, when we aggregate across individuals, balanced growth paths can exist. 
Importantly, balanced growth paths can exist even if a buffer stock target, where $\Ex_{t}[\mNrm_{t+1}/\mNrm_{t}]=1$, does not exist. 
What we require for aggregate stability is the weaker notion of a `pseudo-target' target, namely that there is some ${\mBalLvl}$ such that if $\mNrm_{t}>{\mBalLvl}$, then $\Ex_{t}[\mLvl_{t+1}/\mLvl_{t}] < \PermGroFac$.






\providecommand{\figName}{cNrmTargetFig} 
\providecommand{\figFile}{\figName} %  and on filename
% \figName allows generic definition of hypertargets based on title of figure
\renewcommand{\figName}{cNrmTargetFig} 
\hypertarget{\figFile}{}
\input{Figures/cNrmTargetFig}


%\end{document}\endinput

\hypertarget{onetarget}{}
\hypertarget{Unique-Stable-Points}{}


\subsection{Existence of Target and Pseudo-Target}\label{subsec:onetarget}\hypertarget{TheoremTarget}{}

For both results below, consider the problem defined in Section~\ref{subsec:Setup}.
The first stabiltiy result guarantees the existence of a buffer stock target $\mTrgNrm$ such that if $\mNrm_{t}=\mTrgNrm$, then $\Ex_{t}[\mNrm_{t+1}]=\mNrm_{t}$.
Existence of such a target requires \hyperlink{GICMod}{strong growth impatience}.

\begin{verbatimwrite}{Equations/TheoremMTargetIsStable}
  \begin{theorem}\label{thm:target} (Buffer Stock Target).
% Don't define it if already defined
If \hyperlink{WRIC}{weak return impatience} (Assumption \ref{ass:WRIC}), \hyperlink{FVAC}{finite value of autarky} (Assumption \ref{ass:FVAC}) and \hyperlink{GICMod}{strong growth impatience} (Assumption \ref{ass:GICMod}) hold, then there exists $\mTrgNrm$, with $\mTrgNrm>0$, such that:
    \begin{equation}
      \Ex_t [{\mNrm}_{t+1}/\mNrm_t] = 1 \mbox{~if~} \mNrm_t = \mTrgNrm, 
      \label{eq:mTarget} % Don't define it if already defined
    \end{equation}
    and, 
    \begin{equation}
      \label{eq:stability} % Don't define it if already defined
      \begin{aligned}
        \forall {\mNrm}_t\in(0,\mTrgNrm),      \,\,& \Ex_t [{\mNrm}_{t+1}] > {\mNrm}_t  \\
        \forall {\mNrm}_t\in(\mTrgNrm,\infty), \,\,& \Ex_t [{\mNrm}_{t+1}] < {\mNrm}_t.
      \end{aligned}
    \end{equation}
  \end{theorem}
\end{verbatimwrite} \unskip
\input{Equations/TheoremMTargetIsStable}

\begin{proof}\let\qed\relax
See Appendix \ref{subsubsec:AppxIndividTarget} for the proof.

\end{proof}


%If $\mTrgNrm$ satisfies \eqref{eq:stability}, we say $\mTrgNrm$ is a point of `stability'.
%\end{document}\endinput

\hypertarget{mTargImplicit}{}
Since $\mNrm_{t+1}= (\mNrm_{t}-\cFunc(\mNrm_{t}))\RNrmByGRnd_{t+1}  +\tranShkAll_{t+1}$, the implicit equation for $\mTrgNrm$ becomes:
%
\begin{equation} \label{eq:mTargImplicit}
  \begin{aligned}
    \Ex_{t} [(\mTrgNrm-\cFunc(\mTrgNrm))\RNrmByGRnd_{t+1}+\tranShkAll_{t+1}] & = \mTrgNrm     \\   (\mTrgNrm-\cFunc(\mTrgNrm))\underbrace{\RNrmByG\Ex_{t}[\permShk^{-1}]}_{\colon = \bar{\RNrmByGRnd}}+1 & = \mTrgNrm .
  \end{aligned}
\end{equation}

\hypertarget{Collective-Stability}{}
\hypertarget{pseudo-steady-state}{}


The second, and less restrictive, definition of a target derives from a traditional aggregate question in macro models: whether or not there is a `balanced growth' path in which aggregate variables (income, consumption, market resources) all grow by the same factor $\PermGroFac$.
In particular, if \hyperlink{GICRaw}{growth impatience} holds, the problem will exhibit a `pseudo-target', by which we mean that there is some ${\mBalLvl}$ such that if $\mNrm_{t}>{\mBalLvl}$, then $\Ex_{t}[\mLvl_{t+1}/\mLvl_{t}] < \PermGroFac$.\hypertarget{balgrostable}{}
\hypertarget{balgrostableSolve}{}
Conversely if $\mNrm_{t}<{\mBalLvl}$ then $\Ex_{t}[\mLvl_{t+1}/\mLvl_{t}] > \PermGroFac$.
The pseudo-target $\mBalLvl$ will be such that $\mLvl$ growth matches $\PermGroFac$, allowing us to write the implicit equation for $\mBalLvl$ as follows:
%
\begin{equation}\label{eq:balgrostable}
  \begin{aligned}
    \Ex_{t}[\mLvl_{t+1}]/\mLvl_{t} & =\Ex_{t}[\permLvl_{t+1}]/\permLvl_{t}
    \\  \Ex_{t}[\mNrm_{t+1}\PermGroFac\permShk_{t+1}\permLvl_{t}]/(\mNrm_{t}\permLvl_{t}) & =\Ex_{t}[\permLvl_{t}\PermGroFac\permShk_{t+1}]/\permLvl_{t}
    \\ \Ex_{t}\left[\permShk_{t+1}\underbrace{\left((\mNrm_{t}-\usual{\cFunc}(\mNrm_{t})\Rfree/(\PermGroFac\permShk_{t+1}))+\tranShkAll_{t+1}\right)}_{\mNrm_{t+1}}\right]/\mNrm_{t} & = 1
    \\ 
    \Ex_{t}\left[(\mBalLvl-\usual{\cFunc}(\mBalLvl))\overbrace{\Rfree/\PermGroFac}^{\RNrmByG}+\permShk_{t+1}\tranShkAll_{t+1}\right] & = \mBalLvl
    \\  (\mBalLvl-\usual{\cFunc}(\mBalLvl))\RNrmByG + 1 & = \mBalLvl .
  \end{aligned}
\end{equation}


The only difference between~\eqref{eq:balgrostable} and~\eqref{eq:mTargImplicit} is the substitution of $\RNrmByG$ for $\bar{\RNrmByGRnd}$.\footnote{A third `stable point' is the $\mBalLog$ where $\Ex_{t}[\log \mLvl_{t+1}] = \log \PermGroFac \mLvl_{t}$; this can be conveniently rewritten as $\Ex_{t}\left[\log\left((\mBalLog-\usual{\cFunc}(\mBalLog))\RNrmByGRnd+\permShk_{t+1}\tranShkAll_{t+1}\right)\right]  = \log \mBalLog_{t}$.
Because the expectation of the log of a stochastic variable is less than the log of the expectation, if a solution for $\mBalLog$ exists it will satisfy $\mBalLog > \mBalLvl$; in turn, if $\mTrgNrm$ exists, $\mTrgNrm>\mBalLog$.
The target $\mBalLog$ is guaranteed to exist when the \hyperlink{GICSdl}{log growth impatience} condition is satisfied (see below).
For our purposes, little would be gained by an analysis of this point parallel to those of the other points of stability; but to accommodate potential practical uses,  the {\ARKurl} toolkit computes the value of this point (when it exists) as \texttt{mBalLog}.}$^{,}$
%Theorem~\ref{thm:MSSBalExists}~formally states the relevant proposition.

%\end{comment}

Under the weaker  \hyperlink{GICRaw}{growth impatience} condition, we can verify the existence of this pseudo-target, $\mBalLvl$.

\begin{verbatimwrite}{Equations/TheoremMSSBalExists}
  \begin{theorem}(`Pseudo-Target').
    \labelsafe{thm:MSSBalExists} % Don't define it if already defined
If \hyperlink{WRIC}{weak return impatience} (Assumption \ref{ass:WRIC}), \hyperlink{FVAC}{finite value of autarky} (Assumption \ref{ass:FVAC}) and \hyperlink{GICRaw}{growth impatience} (Assumption \ref{ass:GICRaw}) hold, then there exists a unique $\mBalLvl$, with $\mBalLvl>0$ such that:
    \begin{equation}
      \Ex_t [\permShk_{t+1}{\mNrm}_{t+1}/\mNrm_t] = 1 \mbox{~if~} \mNrm_t = \mBalLvl.
      \label{eq:mBalLvl}
    \end{equation}
    Moreover,
    \begin{equation}
      \label{eq:stabilityStE}
      \begin{aligned}
        \forall {\mNrm}_t\in(0,\mBalLvl),      \,\,& \Ex_{t}[\mLvl_{t+1}]/\mLvl_{t} > \PermGroFac \\
        \forall {\mNrm}_t\in(\mBalLvl,\infty), \,\,& \Ex_{t}[\mLvl_{t+1}]/\mLvl_{t} < \PermGroFac.
      \end{aligned}
    \end{equation}
  \end{theorem}
\end{verbatimwrite} \unskip
\input{Equations/TheoremMSSBalExists}
\begin{proof}
See Appendix \ref{subsubsec:AppxPseudoSS} for the proof.

\end{proof}


\begin{comment}
  \hypertarget{log-pseudo-steady-state}{}
  \subsubsection{The  expected `log-pseudo steady state'}\label{subsubsec:logpseudosteadystate}

  A final point of stability will be indicated by a $\sim$ accent; this is the $\mNrm$ at which the expectation of $\log \mLvl$ is unchanging.

  \begin{verbatimwrite}{Equations/TheoremGroIsStable}
    \begin{theorem}\label{thm:TheoremGroIsStable}{}{\label{thm:TheoremGroIsStable}}
      For the non-degenerate solution to Section~\ref{subsec:Setup}'s problem when {\FVAC}, {\WRIC}, and {\GICRaw} all hold, there exists a unique cash-on-hand-to-permanent-income ratio $\mBalLvl>0$ such that
      \begin{equation}
        \Ex_t [\log {\mLvl}_{t+1} ] = \log \mLvl_t \mbox{~if~} \mLvl_t = \mBalLvl.
        \label{eq:mTarget}
      \end{equation}
      Moreover, $\mBalLvl$ is a point of stability in the sense that
      \begin{equation} \label{eq:stabilityLog}
        \begin{aligned}
          \forall {\mNrm}_t\in(0,\mBalLvl),      \,\,& \Ex_t [\log {\mLvl}_{t+1}] > \log {\mLvl}_t  \\
          \forall {\mNrm}_t\in(\mBalLvl,\infty), \,\,& \Ex_t [\log {\mLvl}_{t+1}] < \log {\mLvl}_t.
        \end{aligned}
      \end{equation}
    \end{theorem}
  \end{verbatimwrite}
  \input{Equations/TheoremGroIsStable}

  with two associated Lemmas demonstrating that $\mBalLvl < \mBalLog$ and $\mBalLvl < \$\mTrgNrm$.
\end{comment}

\subsection{Example With Balanced-Growth \texorpdfstring{$\mBalLvl$}{m} But No Target~\texorpdfstring{$\mTrgNrm$}{m}}

Because the equations defining the buffer stock target and pseudo-target,~\eqref{eq:mTargImplicit} and~\eqref{eq:balgrostable}, differ only by substitution of $\RNrmByG$ for $\bar{\RNrmByGRnd}=\RNrmByG \Ex[\permShk^{-1}]$, if there are no permanent shocks ($\permShk \equiv 1$), the conditions are identical.
For many parameterizations (e.g., under the baseline parameter values used for constructing figure~\ref{fig:cNrmTargetFig}), $\mTrgNrm$ and $\mBalLvl$ will not differ much.

\renewcommand{\figFile}{GICModFailsButGICRawHolds}
\hypertarget{\figFile}{}
\input{Figures/GICModFailsButGICRawHolds}

An illuminating exception is exhibited in Figure~\ref{fig:GICModFailsButGICRawHolds}, which modifies the baseline parameter values by quadrupling the variance of the permanent shocks, enough to cause failure of \hyperlink{GICMod}{strong growth impatience}; now there is no target level of market resources $\mTrgNrm$.
Nonetheless, the pseudo-target still exists because it turns off realizations of the permanent shock.
It is tempting to conclude that the reason target $\mTrgNrm$ does not exist is that the increase in the size of the shocks induces a precautionary motive that increases the consumer's effective patience.
The interpretation is not correct because as market resources approach infinity, precautionary saving against noncapital income risk becomes negligible (as the proportion of consumption financed out of such income approaches zero).
\hypertarget{PermGroFacAdj}{}The correct explanation is more prosaic: The increase in uncertainty boosts the expected uncertainty-modified rate of return factor from $\RNrmByG$ to $\bar{\RNrmByGRnd}>\RNrmByG$ which reflects the fact that in the presence of uncertainty the expectation of the inverse of the growth factor increases: $\PermGroFacAdj < \PermGroFac$.
That is, in the limit as $\mNrm \rightarrow \infty$ the increase in effective impatience reflected in $\GPFacMod > \GPFacRaw$ is entirely due to the elevation of the expected normalized return factor under uncertainty, not to a (limiting) change in precaution.
In fact, the next section will show that an aggregate balanced growth equilibrium will exist even when realizations of the permanent shock are not turned off: The required condition for aggregate balanced growth is the regular \hyperlink{GICRaw}{growth impatience}, which ignores the magnitude of permanent shocks.\footnote{\cite{szeidlInvariant}'s impatience condition, discussed below, also tightens as uncertainty increases, but this is also not a consequence of a precaution-induced increase in patience -- it represents an increase in the tightness of the requirements of the `mixing condition' used in his proof.}



Before we get to the formal arguments, the key insight can be understood by considering an economy that starts, at date $t$, with the entire population at $\mNrm_{t}=\mBalLvl$, but then evolves according to the model's assumed dynamics between $t$ and $t+1$.
Equation~\eqref{eq:balgrostable} will still hold, so for this first period, at least, the economy will exhibit balanced growth: the growth factor for aggregate $\MLvl$ will match the growth factor for permanent income $\PermGroFac$.
It is true that there will be people for whom the financial balances ratio, $\bNrm_{t+1}$, where $\bNrm_{t+1} = \aNrm_{t}\Rfree/(\PermGroFac\permShk_{t+1})$, is boosted by a small draw of $\permShk_{t+1}$.
However, their contribution to the \textit{level} of the aggregate variable is given by $\bLvl_{t+1}=\bNrm_{t+1}\pLvl_{t}\PermGroFac\permShk_{t+1}$, so their $\bNrm_{t+1}$ is reweighted by an amount that exactly unwinds that divisor-boosting.
This means that it is possible for the consumption-to-permanent-income ratio for every consumer to be small enough that their market resources ratio is expected to rise, and yet for the economy as a whole to exhibit a balanced growth equilibrium with a finite aggregate balanced growth steady state $\BalGroFac{\MNrm}$ (this is not numerically the same as the individual \hyperlink{pseudo-steady-state}{pseudo-target} ratio $\mBalLvl$ because the problem's nonlinearities have consequences when aggregated).\footnote{Still, the pseudo-target can be calculated from the policy function without any simulation, and therefore serves as a low-cost starting point for the numerical simulation process; see \href{https://econ-ark.org/materials/harmenberg-aggregation?launch}{Harmenberg-Aggregation} for an example.}


\begin{comment}
 **** candidate for deletion****
\hypertarget{LimitsAsmtToInfty}{}
\subsection{Limits as \texorpdfstring{$\mNrm$}{m} Approaches Infinity}\label{subsec:LimitsAsmtToInfty}

Define
\begin{equation}\begin{gathered}\begin{aligned}
  \Min{\cFunc}(\mNrm)  & = \MPCmin \mNrm
\end{aligned}\end{gathered}\end{equation}
which is the solution to an infinite-horizon problem with no noncapital income ($\tranShkAll_{t+n} = 0~\forall~n\geq1$); clearly $\Min{\cFunc}(\mNrm) < \usual{\cFunc}(\mNrm)$, since allowing the possibility of future noncapital income cannot reduce current consumption.
Our imposition of the {\RIC} guarantees that $\MPCmin > 0$, so this solution satisfies our definition of nondegeneracy, and because this solution is always available it defines a lower bound on both the consumption and value functions.%\footnote{We will assume the \RIC~holds here and subsequently so that $\MPCmin > 0$; the situation is a bit more complex when the \RIC~does not hold.  In that case the bound on consumption is given by the spending that would be undertaken by a consumer who faced binding liquidity constraints.}

Assuming the {\FHWC}~holds, the infinite-horizon perfect foresight solution~\eqref{eq:cFuncPFUnc} constitutes an upper bound on consumption in the presence of uncertainty, since the introduction of uncertainty strictly decreases the level of consumption at any $\mNrm$ (\cite{ckConcavity}).
Thus,
\begin{equation} \label{eq:lowerltupper}
  \begin{aligned}
    \Min{\cFunc}(\mNrm) < & \usual{\cFunc}(\mNrm)  < \bar{\cFunc}(\mNrm)  \\
    1 < & \usual{\cFunc}(\mNrm)/\Min{\cFunc}(\mNrm)  < \bar{\cFunc}(\mNrm)/\Min{\cFunc}(\mNrm).
  \end{aligned}
\end{equation}
But
\begin{equation}\begin{gathered}\begin{aligned}%  \label{eq:limitlowerupper}
  \lim_{m \rightarrow \infty} \bar{\cFunc}(\mNrm)/\Min{\cFunc}(\mNrm)
  & = \lim_{m \rightarrow \infty} (\mNrm -1+ \hNrm)/\mNrm  \\
  & = 1,
\end{aligned}\end{gathered}\end{equation}
so as $\mNrm \rightarrow \infty, \usual{\cFunc}(\mNrm)/\Min{\cFunc}(\mNrm) \rightarrow 1$, and the continuous differentiability and strict concavity of $\usual{\cFunc}(\mNrm)$ therefore implies
\begin{equation*} % \label{eq:limxtoinftycp}
  \lim_{m \rightarrow \infty} \usual{\cFunc}^{\prime}(\mNrm) =
  \Min{\cFunc}^{\prime}(\mNrm) = \bar{\cFunc}^{\prime}(\mNrm) = \MPCmin
\end{equation*}
because any other fixed limit would eventually lead to a level of consumption either exceeding $\bar{\cFunc}(\mNrm)$ or lower than $\Min{\cFunc}(\mNrm)$.
Figure~\ref{fig:mpclimits} illustrates these limits by plotting the numerical solution.


\hypertarget{LimitsAsmtToZero}{}
\subsection{Limits as \texorpdfstring{m}{m} Approaches Zero}
\label{subsec:LimitsAsmtToZero}

Equation~\eqref{eq:MPCmaxInvApndxIter} shows that the limiting value of
$\MPCmax$ is
\begin{equation}\begin{gathered}\begin{aligned}
  \MPCmax  & = 1-{\Rfree}^{-1}{(\pZero  \Rfree\DiscFac)}^{1/\CRRA}.
\end{aligned}\end{gathered}\end{equation}

Defining $\eFunc(\mNrm)=\cFunc(\mNrm)/\mNrm$ we have
\begin{equation}\begin{gathered}\begin{aligned}
  \lim_{m \rightarrow 0} \eFunc(\mNrm)  & = (1-\pZero^{1/\CRRA}\RPFac) = \MPCmax .
\end{aligned}\end{gathered}\end{equation}

Now using the continuous differentiability of the consumption function along with L'H\^opital's rule, we have
\begin{equation}\begin{gathered}\begin{aligned}
  \lim_{m \rightarrow 0} \usual{\cFunc}^{\prime}(\mNrm)  & = \lim_{m \rightarrow 0}
                                                          \eFunc(\mNrm) = \MPCmax.
\end{aligned}\end{gathered}\end{equation}

Figure~\ref{fig:mpclimits} visually confirms that the numerical solution obtains this limit for the MPC as $\mNrm$ approaches zero.
\end{comment}



\end{document}\endinput
